\documentclass[11pt]{article}
\usepackage[margin=1in]{geometry}
\usepackage{amsmath,amssymb,amsthm}
\usepackage[colorlinks=true,linkcolor=blue,urlcolor=blue,citecolor=blue]{hyperref}
\newtheorem{theorem}{Theorem}[section]
\newtheorem{lemma}[theorem]{Lemma}
\newtheorem{corollary}[theorem]{Corollary}
\newcommand{\R}{\mathbb R}
\newcommand{\norm}[1]{\left\lVert#1\right\rVert}
\newcommand{\diag}{\operatorname{diag}}
\newcommand{\spec}{\operatorname{spec}}
\title{Relative spectral universality of one edge insertion\protect\\
in simple unweighted bipartite graphs}
\author{Research working manuscript}
\date{November 2026}
\begin{document}
\maketitle
\begin{abstract}
Fix $0<\alpha<1$. Every finite list $q_i\ge1$ with
$\prod_iq_i=2^\alpha$ is a limit of the nontrivial relative spectra of
fractional graph Laplacians before and after the insertion of one edge.
Both graphs can be simple, unweighted, connected, and bipartite; the
inserted edge has effective resistance exactly one in every host.
All additional relative correction eigenvalues tend to zero.
Consequently the exact worst-case count above energy-relative tolerance
$\varepsilon$ is the largest nonnegative integer strictly less than
$\alpha\log2/\log(1+\varepsilon)$, already in this graph class.
The construction combines a rank-one matrix relative-spectrum theorem,
weighted mass-path realization, bipartite complete-graph mass cells,
and rational resistor gadgets. A compression bound for the fractional
integral avoids any unjustified interchange involving diverging fast
eigenvalues. Historical priority remains under review.
\end{abstract}

\section{Statement and conventions}
For connected graph Laplacians $L,L'$ on the same vertex set, define
the base-energy relative correction
\[
 C_\alpha(L,L')
  =(L^\dagger)^{\alpha/2}(L'^\alpha-L^\alpha)(L^\dagger)^{\alpha/2}.
\]
The Moore--Penrose fractional calculus is zero on constants. The relative
eigenvalues on $\mathbf1^\perp$ are $q_j=1+\lambda_j(C_\alpha)$;
the constant direction supplies a zero correction eigenvalue.
All norms and powers below are Euclidean operator norms and principal
powers. Graphs in the main theorem have ordinary unit vertex masses.

\begin{theorem}[Unweighted bipartite insertion universality]
\label{thm:insertion}
Fix $0<\alpha<1$, a finite list $q_1\ge\cdots\ge q_n\ge1$
with $\prod_iq_i=2^\alpha$, and $\eta>0$.
There exist finite simple unweighted connected bipartite graphs $G,G'$
such that $G'$ is obtained from $G$ by inserting one previously absent
edge $uv$, with
\[
 (e_u-e_v)^T L_G^\dagger(e_u-e_v)=1,
\]
and such that the ordered eigenvalues of $C_\alpha(L_G,L_{G'})$ are
within $\eta$, coordinatewise, of
\[
 (q_1-1,\ldots,q_n-1,0,\ldots,0).
\]
The target list is padded to the constructed graph dimension.
Graph sizes and degrees may grow as $\eta$ decreases.
\end{theorem}

\begin{corollary}[Exact sharp count on this graph class]
\label{cor:count}
For every $\varepsilon>0$, the supremum of
\[
 \#\{j:\lambda_j(C_\alpha(L_G,L_{G'}))>\varepsilon\}
\]
over the graph pairs in Theorem~\ref{thm:insertion} is
\[
 \max\left\{0,\left\lceil
      \frac{\alpha\log2}{\log(1+\varepsilon)}
                         \right\rceil-1\right\}.
\]
In particular, the exact worst-case leading term as
$\varepsilon\downarrow0$ is $\alpha\log2/\varepsilon$.
\end{corollary}

The effective-resistance condition in the theorem is exact at every
finite stage. The result does not assert bounded maximum degree, large
girth, bounded graph size, or exact attainment of a repeated target
relative eigenvalue above one.

\section{Matrix and mass-path input}
We use the following previously proved repository theorem~\cite{Relative}.
For every $\tau>0$, $0<\alpha<1$, and finite admissible list
\[
 q_i\ge1,\qquad \prod_iq_i=(1+\tau)^\alpha,
\]
the list can be approximated arbitrarily closely by the generalized
eigenvalues of $((A+uu^T)^\alpha,A^\alpha)$ with $A\succ0$ and
$u^TA^{-1}u=\tau$. Moreover, after padding the relative correction
by one zero, the same lists are limits of one-edge-conductance increases
on connected weighted paths with fixed positive vertex masses.
The edge increment divided by its base conductance is exactly $\tau$.

For clarity, the matrix construction takes
\[
 r_i=q_i^{1/\alpha},\quad
 z_i=(r_i-1)\prod_{j>i}r_j,\quad
 a_i=\delta^{n-i},\quad A=\diag(a_i),\quad u_i=\sqrt{a_i z_i}.
\]
The relative correction tends to $\diag(q_i-1)$ while its leverage is
exactly $\sum_i z_i=\tau$. The mass-path transfer is by orthogonal
Lanczos and the grounded Jacobi construction, with an additional
ground vertex of mass tending to infinity. These input results are
proved in full in~\cite{Relative}; they do not require a solver or
finite-precision assertion.

It remains to simulate a fixed weighted mass path by simple unweighted
bipartite graphs while controlling the fractional relative spectrum.
This additional step is proved below.

\section{A fractional compression lemma}
\begin{lemma}[Converging compressed powers]\label{lem:compression}
Let $A_s\succeq0$ act on finite-dimensional spaces whose dimensions
may grow with $s$. Let $J_s:\R^d\to\R^{N_s}$ be isometries, where
$d$ is fixed. Suppose that for every fixed $t>0$,
\[
 J_s^T(A_s+tI)^{-1}J_s\longrightarrow(H+tI)^{-1}
\]
for a fixed $H\succeq0$, and suppose
$\sup_s\norm{J_s^TA_sJ_s}<\infty$. Then for every $0<\alpha<1$,
\[
 J_s^T A_s^\alpha J_s\longrightarrow H^\alpha
 \quad\hbox{in operator norm}.
\]
\end{lemma}
\begin{proof}
Let $M$ bound $\norm{J_s^TA_sJ_s}$. The compression inverse inequality
\[
 J_s^T(A_s+tI)^{-1}J_s
       \succeq(J_s^TA_sJ_s+tI)^{-1}
\]
follows from the Schur-complement inverse in an orthogonal basis
beginning with $J_s$. Consequently
\begin{align*}
 0&\preceq J_s^T A_s(A_s+tI)^{-1}J_s\\
  &\preceq (J_s^TA_sJ_s)(J_s^TA_sJ_s+tI)^{-1}
   \preceq\min(1,M/t)I.
\end{align*}
The last bound is integrable against $t^{\alpha-1}dt$ on
$(0,\infty)$ precisely for $0<\alpha<1$.
The fractional representation
\[
 A^\alpha=\frac{\sin(\pi\alpha)}\pi
             \int_0^\infty t^{\alpha-1}A(A+tI)^{-1}\,dt
\]
and dominated convergence now prove the assertion, since the assumed
resolvent limit makes the compressed integrand tend to
$H(H+tI)^{-1}$.
\end{proof}

This lemma is why no uncontrolled estimate of a growing high-eigenvalue
fractional power enters the proof.

\section{Bipartite cells and resistor gadgets}
Fix a connected conductance path on $d$ macro vertices, with fixed
positive masses $M_i$. Its conductance Laplacian is $K$, and its
mass-normalized Laplacian is $H=M^{-1/2}KM^{-1/2}$.
Choose a distinguished path edge. In the bridge-doubling version,
normalize its conductance to one by a common scaling of $K$ and $K'$.
Their relative spectrum is invariant under this common scaling.
Other conductances can be approximated by positive rational numbers,
without changing the distinguished edge. Continuous fixed-rank
functional calculus preserves the desired relative spectral approximation.
Hence assume from now on that all nondistinguished conductances are
positive rationals.

For a rational macro conductance $p/q>0$, an unweighted network of
$2p$ internally disjoint paths, each of length $2q$, has effective
conductance $p/q$ between its two terminals. The terminals have the
same bipartite color, and the network is simple, including when $q=1$:
the paths then have length two and distinct internal vertices.

As integers $s$ tend to infinity, replace vertex $i$ by the
unweighted complete bipartite cell $K_{m_i(s),m_i(s)}$, where
\[
 m_i(s)\ge1,\qquad \frac{2m_i(s)}s\longrightarrow M_i.
\]
Choose one port vertex in each cell. Every external connection to
that cell, including all ends of all resistor gadgets, uses this
same single port. Replace each nondistinguished macro edge by its
fixed unweighted even-path gadget between the corresponding ports.
Cells on the same side of the distinguished edge use ports of the
same global color; across that edge flip the chosen port color.
The actual distinguished edge is a single unit-conductance port edge.
All gadget internal vertices are new and disjoint from each other and
from the cells.

The resulting host is simple, connected, unweighted, and bipartite.
The distinguished edge is a bridge. Increasing its conductance from
one to $1+\tau$ has exact leverage $\tau$ for every $s$.
Cell sizes and port degrees are unrestricted.

\begin{lemma}[Low compressed resolvent]\label{lem:resolvent}
Let $L_s$ be the host Laplacian in this construction, and put
$A_s=sL_s$. Let $J_s$ have as its $i$th column the normalized indicator
of cell $i$, and be zero at gadget vertices. Then for every fixed $t>0$,
\[
 J_s^T(A_s+tI)^{-1}J_s\longrightarrow(H+tI)^{-1},
 \qquad \sup_s\norm{J_s^TA_sJ_s}<\infty.
\]
The same conclusion holds for the updated graph and its updated
macro matrix $H'$.
\end{lemma}
\begin{proof}
Decompose the physical space orthogonally into the cell-constant
space $P=\operatorname{range}J_s$, the gadget-coordinate space $Z$,
and the mean-zero cell space $F$. The dimension of $P$ and of $Z$
is fixed; the dimension of $F$ may grow. Write
\[
 R_s=\diag\sqrt{\frac{s}{2m_i(s)}}\longrightarrow M^{-1/2}.
\]
Collapse each cell to one terminal in the fixed external gadget
network. Its conductance matrix on terminal and gadget vertices is
\[
 K_{\rm ext}=\begin{pmatrix}K_{PP}&K_{PZ}\\K_{ZP}&K_{ZZ}\end{pmatrix}.
\]
Here $K_{ZZ}\succ0$ when gadgets exist: each gadget component has a
Dirichlet boundary in the terminal set. Its Schur complement is exactly
the prescribed macro path Laplacian,
\[
 K=K_{PP}-K_{PZ}K_{ZZ}^{-1}K_{ZP}.
\]
If there are no gadget vertices, omit the $Z$ block and set $K=K_{PP}$.

Before eliminating $F$, the blocks of $A_s+tI$ satisfy
\[
 \begin{array}{lll}
 PP=R_sK_{PP}R_s+tI,&
 PZ=\sqrt{s}R_sK_{PZ},& ZZ=sK_{ZZ}+tI,\\
 PF=O(\sqrt{s}),& ZF=O(s),& FF\succeq c s^2 I.
 \end{array}
\]
All norm constants are independent of $s$. Indeed the number of
external edges and ports is fixed, and the normalized cell-indicator
values at a port are $O(s^{-1/2})$. The unscaled cell Laplacian has
nonzero eigenvalues $m_i(s)$ and $2m_i(s)$, so after multiplication
by $s$ its mean-zero block is bounded below by $c s^2I$.
The external Laplacian is PSD and cannot reduce this bound.

Eliminating $F$ therefore changes the $PP$ block by $O(s^{-1})$,
the $PZ$ block by $O(s^{-1/2})$, and the $ZZ$ block by $O(1)$.
Eliminating $Z$ next gives the Schur complement on $P$,
\[
 R_s(K_{PP}-K_{PZ}K_{ZZ}^{-1}K_{ZP})R_s+tI+o(1)
       \longrightarrow H+tI.
\]
Its inverse is the compressed resolvent in the statement.
The explicit $PP$ block above also proves the uniform compression
bound. Updated conductances change only fixed external blocks, so
the same argument applies to $L_s'$.
\end{proof}

\section{Ritz lower bounds and exhaustion of the determinant budget}
\begin{lemma}[Full relative spectral transfer]\label{lem:transfer}
In the bridge construction above, the relative correction eigenvalues
of the ordinary graph pair converge, with extra zero entries, to
those of the fixed weighted mass-path pair.
\end{lemma}
\begin{proof}
Lemmas~\ref{lem:resolvent} and~\ref{lem:compression} give
\[
 J_s^T(sL_s)^\alpha J_s\to H^\alpha,\qquad
 J_s^T(sL_s')^\alpha J_s\to H'^\alpha.
\]
Let $g_s=(\sqrt{2m_i(s)})_{i=1}^d$, and choose isometries
$V_s:\R^{d-1}\to g_s^\perp$ converging to an isometry
$V:\R^{d-1}\to(\sqrt{M_i})^\perp$.
Every vector in $\operatorname{range}(J_sV_s)$ is orthogonal to the
physical constant vector, because it is zero at gadget vertices and
its cell-coordinate sum is $g_s^TV_s=0$.
The compressed positive-definite pair
\[
 \left(V_s^TJ_s^T(sL_s')^\alpha J_sV_s,
       V_s^TJ_s^T(sL_s)^\alpha J_sV_s\right)
\]
converges to the corresponding positive-definite pair for
$(H'^\alpha,H^\alpha)$ on their common positive support.
Its ordered generalized Ritz eigenvalues $\mu_{s,1},\ldots,
\mu_{s,d-1}$ thus converge to the macro relative eigenvalues
$q_1,\ldots,q_{d-1}$.

Let $\lambda_{s,1}\ge\cdots\ge\lambda_{s,N_s-1}\ge1$
be the full relative eigenvalues on $\mathbf1^\perp$.
The generalized min--max principle yields
\[
 \lambda_{s,j}\ge\mu_{s,j}\ge1,\qquad 1\le j\le d-1.
\]
The bridge leverage is exactly $\tau$ in both the physical and the
macro graph. The determinant lemma on their common positive supports
therefore gives
\[
 \sum_{j=1}^{N_s-1}\log\lambda_{s,j}
    =\alpha\log(1+\tau)
    =\sum_{j=1}^{d-1}\log q_j.
\]
As $s\to\infty$, the sum of the logarithms of the Ritz eigenvalues
tends to this whole budget. Hence
\begin{align*}
 &\sum_{j=1}^{d-1}(\log\lambda_{s,j}-\log\mu_{s,j})
       +\sum_{j=d}^{N_s-1}\log\lambda_{s,j}\\
 &\hspace{25mm}
   =\alpha\log(1+\tau)-\sum_{j=1}^{d-1}\log\mu_{s,j}
       \longrightarrow0.
\end{align*}
Every term on the left is nonnegative. Thus the first $d-1$ full
relative eigenvalues converge individually to $q_j$, and the largest
remaining relative eigenvalue tends to one. Subtracting one, and
including the constant correction eigenvalue zero, proves the claim.
The common time scaling by $s$ cancels from every relative matrix.
\end{proof}

\section{Making the update an actual unweighted edge insertion}
For $\tau=1$, replace the distinguished unit bridge by three internally
disjoint paths of length three between its two ports. This fixed
unweighted simple bipartite gadget has effective conductance one,
and its terminals have opposite global colors. There is no direct
edge between them. Insert that direct unit edge for the update.
After Schur elimination, the two macro conductances are one and two,
as required. The construction in Lemma~\ref{lem:resolvent} applies
unchanged, adding the finitely many new gadget vertices to $Z$.

The resistance between these ports is \emph{exactly one} in every
finite host. All external connections of a macro cell use the same
port, so its complete bipartite cell hangs from a single terminal.
Each component of the rest of the macro path also attaches only at
one of these two terminals. Such dangling networks carry no current
in the unit terminal-to-terminal electrical problem. The three
length-three paths are therefore the entire effective connector and
have resistance $3/3=1$.

Thus the inserted physical edge has exact leverage one, and the
updated graph is again simple, unweighted, connected, and bipartite.
The determinant budget in Lemma~\ref{lem:transfer} is still exactly
$\alpha\log2$ at every finite stage. Its proof and conclusion hold
for this insertion construction as well.

\begin{proof}[Proof of Theorem~\ref{thm:insertion}]
Use the matrix and mass-path input with $\tau=1$ to approximate the
prescribed list. Commonly scale the macro conductances so the edited
edge has base conductance one and increment one. Approximate the
other macro conductances by positive rationals. Realize the resulting
mass path by the bipartite cell and resistor construction, using the
three-path insertion gadget at the distinguished edge.
Lemma~\ref{lem:transfer} gives relative spectral approximation with
all extra correction eigenvalues tending to zero. Choose sufficiently
large finite parameters at each of these sequential approximation
steps, allocating the prescribed tolerance among them. The inserted
edge has exactly unit effective resistance throughout the final
cell-size sequence, by the preceding electrical argument.
\end{proof}

\begin{proof}[Proof of Corollary~\ref{cor:count}]
For any connected host with a unit edge insertion of effective
resistance one, operator monotonicity gives $\lambda_j(C_\alpha)\ge0$
and the determinant lemma gives
\[
 \prod_j(1+\lambda_j(C_\alpha))=2^\alpha.
\]
If $m$ eigenvalues exceed $\varepsilon$, then
$(1+\varepsilon)^m<2^\alpha$. The stated largest strict integer is
therefore an upper bound, including at integral thresholds.
For the matching lower bound, if that integer is $m>0$, take the
target list of $m$ entries equal to $2^{\alpha/m}>1+\varepsilon$
and use Theorem~\ref{thm:insertion} with a sufficiently small strict
margin. If $m=0$, the upper bound is already exact.
\end{proof}

The same proof also gives the following additional class, which keeps
an exact arbitrary leverage parameter but allows one weighted edge
after updating.
\begin{corollary}[Arbitrary bridge increase]\label{cor:bridge}
For every $\tau>0$, every finite list $q_i\ge1$ with
$\prod_iq_i=(1+\tau)^\alpha$ is approximable, with extra relative
correction eigenvalues tending to zero, by a simple unweighted
connected bipartite host whose single bridge edge is increased from
weight one to $1+\tau$.
The exact strict-count supremum on this class is
\[
 \max\left\{0,\left\lceil
    \frac{\alpha\log(1+\tau)}{\log(1+\varepsilon)}
                           \right\rceil-1\right\}.
\]
The host is unweighted; for general $\tau$ the updated graph has one
edge of weight $1+\tau$.
\end{corollary}

\section{Scope and priority}
The theorem is about energy-relative fractional corrections and permits
unbounded graph size and degree. It gives no runtime theorem, no
bounded-degree construction, and no finite-word condition bound.
The exact count grows as $1/\varepsilon$, unlike the condition-free
logarithmic-squared \emph{absolute} compression bounds in the
repository. These two error criteria differ materially.

Resistor series and parallel laws, Schur complements, complete-bipartite
spectra, the compression inverse inequality, generalized Ritz
interlacing and the determinant identity are classical tools. The
coarse-graining mechanism of resolvent convergence under strong
intra-cluster coupling is also established; see Koke~\cite{Koke},
especially the undirected results in Section 3. Here the particular
growing unweighted cells and fixed resistor gadgets are treated
directly, and the relative spectrum is controlled by the additional
fractional compression and determinant-budget arguments.
The
rank-one relative spectral closure and weighted mass-path transfer
are separate inputs proved in~\cite{Relative}.
The present possible contribution is their sharp realization by
simple unweighted bipartite graph pairs differing by one inserted
edge of exactly fixed effective resistance. Historical novelty
remains unverified; this source does not certify a seminal advance.

\begin{thebibliography}{9}
\bibitem{Relative} Every determinant-admissible relative spectrum can arise
from one edge change. Repository working manuscript,
\texttt{spectral-gap/relative-fractional-response.tex}.
\bibitem{BCKS} B. Beckermann, A. Cortinovis, D. Kressner, M. Schweitzer.
Low-rank updates of matrix functions II: Rational Krylov methods.
SIAM J. Numer. Anal. 59 (2021), 1325--1347.
\url{https://arxiv.org/abs/2008.11501}.
\bibitem{Koke} C. Koke.
Di-graphs with tightly connected clusters: Effective Graph Laplacians
and resolvent convergence. arXiv:2601.18057v2 (2026).
\url{https://arxiv.org/abs/2601.18057}.
\end{thebibliography}
\end{document}
